\documentclass[11pt]{article}
\usepackage[margin=1in]{geometry}
\usepackage{amsmath,amssymb,amsthm,hyperref}
\newtheorem{proposition}{Proposition}
\newtheorem{lemma}{Lemma}
\title{Hillclimb 003: exact rank, factored joins, and bounded-source C38 tensors}
\author{Structural research iteration; no new ECDLP runtime bound}
\date{31 October 2026}
\begin{document}
\maketitle

\paragraph{Status.}
The goal remains a correct classical ECDLP algorithm with expected bit
time $2^{\widetilde O(\sqrt b)}$ on the stated general prime-field class.
The inherited contraction and reduction theorems remain under audit.
This iteration supplies exact representation counterexamples, a positive
factored-join example, and an executable whole-graph refactoring that
removes the arity factor from the root degree bound. It does not prove
closure of compact representations through the full recursion or an
improved ECDLP runtime. Finite checks and algebraic arguments are identified
separately below.

\section{Large Fourier-cut rank in the actual first predicate}
Write $B=2^k$ for the arithmetic limb base, $L=L_1L_2$ for the
transform grid, and $M_F$ for its Fourier modulus. At first vertex $F_0$,
a source word with raw coefficients $v_0,\ldots,v_{d-1}$ at positions
$0,L_2,\ldots,(d-1)L_2$, and zero elsewhere, has Fourier row
\[
 h_j=\sum_{s=0}^{d-1}v_s\omega_1^{sj}\pmod{M_F},\qquad
 \omega_1=\omega^{L_2}.
\]
Choose $d\le L_1/2$ and use quotient sources $q_i$. Their local range
tests do not interact with the distinctness tests on point coordinates.

\begin{lemma}[A unit identity minor]
Suppose every chosen digit $v_s$ ranges independently over a set $A_s$
while all range-state ports remain fixed, and all other source ports are
fixed to a valid local assignment. For $h$ independently varied quotient
sources, the first predicate's matrix across the cut after the first $d$
Fourier words has an identity minor of dimension
\[
 N=\left(\prod_{s=0}^{d-1}|A_s|\right)^h.
\]
The statement remains true after canonical recoding of the complete words
into Boolean or simplex symbols; the cut groups complete words.
\end{lemma}
\begin{proof}
Any $d$ consecutive Fourier evaluations determine the $d$ raw coefficients:
the associated Vandermonde determinant is a product of root differences
and units. All these differences are units in the supplied transform
ring. Select a row index from the first $d$ evaluations of one coefficient
vector and a column index from the remaining evaluations of another.
The latter include another $d$ consecutive evaluations. If the assembled
row is valid, invertibility forces both coefficient vectors to be equal.
When they are equal, all pinned range states and other factors accept.
This is precisely an identity submatrix. Different quotient sources have
disjoint ports, so the dimensions multiply.
\end{proof}

For completeness, common range pins can be chosen constructively. Scan
the base-$B$ digits of the bound $16p$, keeping the borrow bit fixed at
each step. At a free digit the two possible borrow outcomes split
$\{0,\ldots,B-1\}$ into two intervals; the larger has at least $B/2$
elements. Use that interval and its common outgoing state. Fix all other
digits to zero, and place the free digits strictly below the highest
nonzero bound digit. The highest digit then forces the final borrow to
one, so the upper-bound test accepts. This gives $|A_s|\ge B/2$.

In the intended grid regime $k=\Theta(\log^2 b)$ and
$L=\Theta(b/\log^2b)$, one can use
$d=\Theta(\sqrt b/\log b)$ consecutive local coefficients, and there
are $h=\Theta(m)$ independent quotient sources. Hence
\[
 \log_2N\ge hd(k-1)=\Omega(m\sqrt b\log b).
\]
This rules out small exact rank across this \emph{specified Fourier cut}.
It is not a lower bound for every tensor-train ordering, change of basis,
arithmetic circuit, or contraction algorithm. An ordering interleaving
independent components can have far smaller ranks.

The checker evaluates six actual minors using two choices per digit,
through dimension 64, and verifies the transform determinants are units.
The larger dimensions from all allowed digits are consequences of the
lemma, not matrices that the program enumerated.

\section{Actual projected entries: high rank and a short formula coexist}
Use the finite general-prime fixture $p=13,L=8$, with $L_1=2$ and
$k=4$, in the Boolean alphabet. For each selected quotient source, its
two Fourier words at $F_0$ are equal to its single nonzero raw limb.
Pin the high limb bit to one, the higher Fourier bits to zero, and range
states to their common values. Each lower free bit becomes an equality
pair. Selecting $r$ such pairs gives the actual restricted C38 tensor
\[
 T(y,z)=\prod_{j=1}^{r}\mathbf1_{y_j=z_j}.
\]
This is a restriction of one local tensor, not of the complete relation
system. The other endpoints in this probe are explicitly all-one tensors.

Give the C38 tensor label $(0,0)$ and its two sets of endpoints labels
$(1,1)$ and $(2,2)$. With packet parameters $a=u=2$, the executable
packetizer puts each real edge at position zero in its own packet. At
node $\tau=1$, there is one surviving mode per packet, its source form
is the bit itself, and its row coefficient is the same polynomial $c(Z)$.
It has $c(0)=0$ and $v_2([Z]c)=3$.

\begin{proposition}[Factored projected table]
For binary mode-exponent vectors $\alpha,\beta\in\{0,1\}^r$, the
target group's exact projected table is
\[
 G(\alpha,\beta)=\prod_{j=1}^{r}K_{\alpha_j,\beta_j},\qquad
 K=\begin{pmatrix}1&-c\\-c&2c^2\end{pmatrix}.
\]
All entries with any exponent greater than one vanish. With the actual
all-one source tensors included, the complete join equals
\[
 (4-4c+2c^2)^r.
\]
\end{proposition}
\begin{proof}
Sum the two allowed assignments of one equality pair. Its generating
polynomial is
$(1-cz)(1-cw)+c^2zw=1-cz-cw+2c^2zw$, giving $K$.
Different pairs multiply. A source bit has moment two at exponent zero
and one at exponent one. Thus the joined factor is
$4K_{00}+2K_{01}+2K_{10}+K_{11}=4-4c+2c^2$.
\end{proof}

The two exterior source tables must be charged and included. Summing
only $G$ would give $(1-2c+2c^2)^r$, whose constant term is one;
the full join's constant term is $4^r$. This erroneous shortcut was
detected and corrected before the final run.

\begin{proposition}[A finite-precision obstruction to block low rank]
Let $r$ be even. Work in
$A=(\mathbb Z/2^P\mathbb Z)[Z]/(Z^q)$ with $P>3r$, $q>r$.
Restrict both exponent vectors to Hamming weight $r/2$.
Any exact separated representation of that block as
$\sum_{i=1}^{\chi} a_i(\alpha)b_i(\beta)$ over $A$ satisfies
\[
 \chi\ge\frac{\binom r{r/2}}{r+1}.
\]
\end{proposition}
\begin{proof}
Write $c_1=[Z]c=2^3v$ with $v$ odd. The coefficient of $Z^r$ in
the block is
\[
 C_{U,V}=c_1^r2^{|U\cap V|},\qquad |U|=|V|=r/2.
\]
After removing the common factor $2^{3r}$, reduction modulo two is
the permutation matrix for complementation: the intersection is empty
exactly when $V$ is the complement of $U$. Hence over
$R=\mathbb Z/2^P\mathbb Z$ the matrix $C$ has
$\binom r{r/2}$ nonzero Smith factors, all of valuation $3r$.
Its inner rank is that number. One way to see this is to diagonalize
the left factor of a putative product $AB$: with inner dimension $s$,
invertible row operations make at most $s$ rows of $AB$ nonzero, whereas
the row module of $C$ needs one generator per nonzero Smith factor.
The latter follows by reducing that row module modulo its double.

Extracting the coefficient of $Z^r$ from each product
$a_i(\alpha)b_i(\beta)$ gives at most $r+1$ outer products over $R$.
Thus the inner rank of $C$ is at most $\chi(r+1)$, proving the bound.
The chosen indices have total degree $r$ and mode weight $3r$, so they
survive the specified jet and precision cutoffs.
\end{proof}

Despite this rank, the formula for $G$ has $O(r)$ factors, and its
complete join above is computable with polynomial coefficient overhead.
This is a positive example of avoiding projected-table enumeration.
It does not establish that the full C38 computation has this form.
These singleton packets are a valid finite geometry but do not obey the
asymptotic balanced packing schedule. No claim about every scheduled
C38 projected table follows from this example.

The program checks 340 target projected entries for $r=1,2,3,4$,
680 actual exterior child contractions, protected rank slices, and
independent full-deformation sums for $r\le2$. Larger $r$ values in
the record are algebraic family certificates, not executed full tables.

\section{A whole-network refactoring with paid copy wires}
The factored example suggests preserving arithmetic factors instead of
merging them prematurely. The actual general-prime C38 network admits
the following exact refactoring.

\begin{proposition}[Bounded-source C38 compiler]
Let $S_F$ be the old number of forward Fourier-symbol edges and $S$
the total edge count. There is a parsimonious graph refactoring with
\[
 S'=S+LW\sum_{s:d_s>1}d_s\le S+6S_F\le7S,
\]
where $W$ is the canonical Fourier word width in the chosen alphabet
and $d_s$ is the number of output equations consuming source $s$.
Every new first-stage arithmetic factor uses at most two source arrays;
every middle equation uses at most eight source arrays; every closure
factor uses one output array. Copy tensors have degree at most seven.
\end{proposition}
\begin{proof}
The first predicate splits into disjoint port sets. For each $j\ge2$,
pair $x_j$ with $x_1$ if $j=2$, otherwise with $X_{j-1}$; include
both range tests and the one distinctness state chain. These pairs use
each point/partial-abscissa source exactly once. All remaining ordinate,
slope and quotient sources have independent range tests and form singleton
factors. No first-stage port is copied. The closure predicate is a product
over equations with disjoint Fourier and carry ports, so it also splits
without new wires.

For the middle stage, the square/chord equations use at most six sources,
and the alternating ordinate-link equation uses at most eight. A source
$x_j$ occurs in at most its two chord equations and its hash; a slope
occurs in its two chord equations and at most one link; a partial abscissa
occurs in four neighboring chord equations and at most two links; an
even ordinate occurs in at most two chord and two link equations; a
quotient occurs only in its own equation. Thus $d_s\le6$.

For a source used by one equation, redirect its existing Fourier edge
to that equation tensor. Otherwise retain the old edge to an equality
tensor and add $d_s$ edges, one to each consuming equation tensor.
Do this for every symbol in every forward word. The equation tensor
performs the same transform, arithmetic operation, and inverse-transform
check for its one output. All its omitted input sources are outside
that equation's syntactic dependency set.

Each original accepting assignment extends uniquely: every copy branch
must equal its old source edge. Conversely, every accepted refactored
assignment restricts to an original one. This proves both parsimony and
the exact added-wire ledger. Input programs, fixed templates, transforms
and graph construction all have polynomial cost in $S$ and $m$.
\end{proof}

The current implementation uses the original first predicate with valid
fixed templates for the omitted disjoint factors. The templates are
constructed from public bounds; they do not contain a planted relation.
This can duplicate polynomial description work, which is explicitly
allowed and charged. Middle and closure equation descriptors read only
their incident ports. Noncanonical words and nonbinary side states are
still rejected.

\paragraph{Degree accounting.}
First-factor degree is at most
$2L_1W+6\lceil n/L_2\rceil$ up to endpoint constants.
Middle-equation degree is at most $9L_2W$.
Closure-equation degree is at most
$RL_1W+2\lceil L/L_2\rceil W_c$, with carry width $W_c$.
Other tensors have constant degree. Parallel multiplicity is at most
$\max\{RW,W,1\}$. These bounds have no multiplicative $m$ factor.
In the intended Fourier regime they give
\[
 \Delta'=\widetilde O(\sqrt b),\quad
 \mu'=\operatorname{polylog}b,\quad
 S'=\widetilde O(mb).
\]
Thus the specific root-degree obstruction identified in iteration 002
is removed even for polynomially growing $m$. This does not certify
the entire relation/reduction proof at such arities, nor change the
current contractor's $S/\sqrt{\log S}$ cost scale. Added wires can
worsen its leading constant. No coefficient improvement is claimed.

Four full finite refactorings pass 324 local-product comparisons, 48
whole-graph mutation comparisons and 32 inconsistent-copy controls.
The dependency program is also scanned for all arities $2$ through 100;
the all-arity bounds come from the source-occurrence argument above.
In the $m=8$ fixture the maximum degree changes from 4,013 to 288,
while the edge count changes from 13,754 to 24,890. No full unpinned
recursive contraction of this refactored graph has been performed.

\section{The remaining closure problem}
A small arithmetic circuit is not by itself a fast contraction algorithm.
For example, the product of the multilinear Boolean clause predicates
of a width-three CNF has a small factored description; summing that
product on Boolean inputs counts satisfying assignments. This observation
is a reduction, not a lower-bound theorem, but prevents treating circuit
description size as a paid summation bound.

The next hypothesis must therefore specify a structural invariant:
which source/mode scopes occur, how they change through projection and
coefficient extraction, and how much state is required for an exact join.
Acceptance requires a construction/access/composition bound for the actual
refactored C38 family through all recursion levels and a complete ECDLP
arity balance. Counterexamples should target that proposed invariant.
The positive copy-factor family above establishes only one solvable case.
The active goal remains unachieved.
\end{document}
